\clearpage\begingroup
\SourceLeaf{136}\NoSourceHead
\newcommand{\IVpoint}{\mathpunct{\vcenter{\hbox{.}}}}
\begin{center}
[IV.]\par\bigskip
\textbf{AUFGABE [AUS DER DIOPTRIK].}\par\medskip
\rule{.40\linewidth}{.4pt}\par
[Aus Handbuch 19, Be, S. 90.]\par
\rule{.40\linewidth}{.4pt}
\end{center}
\noindent\hspace*{\parindent}Ein Lichtstral fällt auf eine Kugelfläche, deren Halbmesser $=r$, Brechungsverhältniss $n:1$ ist. Abscissen $x$ werden vom Centrum gezählt. Vor der Brechung sind die Gleichungen der Lichtbahn
\[
\begin{aligned}
y&=\alpha+\mu x\\
z&=\beta+\nu x,
\end{aligned}
\]
\noindent nach der Brechung hingegen
\[
\begin{aligned}
y&=\alpha'+\mu'x\\
z&=\beta'+\nu'x.
\end{aligned}
\]
\noindent Man soll $\alpha',\beta',\mu',\nu'$ aus $\alpha,\beta,\mu,\nu$ bestimmen.

Ist $\mu=0$, $\nu=0$, so trifft der gebrochene Stral die Axe in
\[
x=\frac{\surd(nnrr-\alpha\alpha-\beta\beta)+\surd(rr-\alpha\alpha-\beta\beta)}{nn-1}
\]
\noindent und man hat
\[
\begin{aligned}
\alpha'&=\alpha\mathbin{.}\frac{\surd(nnrr-\alpha\alpha-\beta\beta)+\surd(rr-\alpha\alpha-\beta\beta)}{\surd(nnrr-\alpha\alpha-\beta\beta)+nn\surd(rr-\alpha\alpha-\beta\beta)}\\[6pt]
\mu'&=-\frac{\alpha(nn-1)}{\surd(nnrr-\alpha\alpha-\beta\beta)+nn\surd(rr-\alpha\alpha-\beta\beta)}\IVpoint
\end{aligned}
\]

Ist $\beta=0$, $\nu=0$, so hat man,
\[
\begin{aligned}
\surd((1+\mu\mu)rr-\alpha\alpha)&=P\\
\surd((1+\mu\mu)nnrr-\alpha\alpha)&=Q\\
nnP+Q+\alpha\mu(nn-1)&=N
\end{aligned}
\]
\noindent gesetzt,
\par\HistoricalFoot{}{17\textsuperscript{*}}
\clearpage\SourceLeaf{137}\HistoricalHead{132}{L}{PHYSIK.\quad NACHLASS.}
\[
\begin{aligned}
\alpha'&=\frac{\alpha(P+Q)}N=\alpha-\frac{\alpha(nn-1)(P+\alpha\mu)}N\\[6pt]
\mu'&=\mu-\frac{\alpha(1+\mu\mu)(nn-1)}N\IVpoint
\end{aligned}
\]
\noindent\hspace*{\parindent}Macht man
\[
\alpha=\lambda\surd(1+\mu\mu),\qquad \alpha'=\lambda'\surd(1+\mu'\mu'),
\]
\noindent so wird
\[
\begin{aligned}
\lambda'&=\frac\lambda n\\[6pt]
\mu'&=\mu-\frac{\lambda(1+\mu\mu)(nn-1)}{nn\surd(rr-\lambda\lambda)+\surd(nnrr-\lambda\lambda)+(nn-1)\lambda\mu}\IVpoint
\end{aligned}
\]
\begin{center}\rule{.25\linewidth}{.4pt}\end{center}
\clearpage\endgroup
